\begin{figure}[ht]
\centering
\begin{tikzpicture}[x=1cm, y=1cm, font=\small]

% Mean value theorem: secant from (a,f(a)) to (b,f(b)) is parallel to
% the tangent at some interior c. f(x) = 0.5 + 1.6*sin(0.5*x) on [0.5, 7].

\def\f(#1){(2.0 + 1.3*sin(deg(0.45*(#1))))}

% Axes
\draw[->, thick] (-0.3, 0) -- (8.2, 0) node[right] {$x$};
\draw[->, thick] (0, 0) -- (0, 4.4) node[above] {$y = f(x)$};

% Curve
\draw[thick, engblue] plot[domain=0.4:7.6, samples=100] (\x, {\f(\x)});

% Endpoints a = 1, b = 7
\def\xa{1.0}\def\xb{7.0}
\pgfmathsetmacro{\fa}{2.0 + 1.3*sin(deg(0.45*\xa))}
\pgfmathsetmacro{\fb}{2.0 + 1.3*sin(deg(0.45*\xb))}
\fill[engred] (\xa, {\fa}) circle (1.8pt);
\fill[engred] (\xb, {\fb}) circle (1.8pt);
\node[anchor=north, font=\scriptsize] at (\xa, -0.05) {$a$};
\node[anchor=north, font=\scriptsize] at (\xb, -0.05) {$b$};

% Secant
\pgfmathsetmacro{\slope}{(\fb - \fa)/(\xb - \xa)}
\draw[thick, engamber] (\xa, {\fa}) -- (\xb, {\fb});
\node[engamber, anchor=south, font=\scriptsize] at (4.0, {\fa + \slope*3.0 + 0.1}) {secant, slope $\frac{f(b)-f(a)}{b-a}$};

% c where f'(c) = slope. f'(x) = 1.3*0.45*cos(0.45 x) = slope =>
% cos(0.45 c) = slope/(0.585). We solve numerically: pick c near 2.2.
% slope = (fb-fa)/6. For these values f'(c)=slope has a solution; mark c approx 2.15.
\def\xc{2.15}
\pgfmathsetmacro{\fc}{2.0 + 1.3*sin(deg(0.45*\xc))}
\fill[englime] (\xc, {\fc}) circle (1.8pt);
\draw[dashed, enggray] (\xc, 0) -- (\xc, {\fc});
\node[anchor=north, font=\scriptsize] at (\xc, -0.05) {$c$};
% tangent at c parallel to secant
\draw[thick, englime] (\xc-1.2, {\fc - \slope*1.2}) -- (\xc+1.2, {\fc + \slope*1.2});
\node[englime, anchor=west, font=\scriptsize] at (\xc+1.0, {\fc + \slope*1.0 + 0.2}) {tangent at $c$ (parallel)};

\end{tikzpicture}
\caption[The mean value theorem]{For $f$ continuous on $[a,b]$ and
differentiable on $(a,b)$, there is an interior point $c$ where the
instantaneous rate of change $f'(c)$ equals the average rate of
change $(f(b)-f(a))/(b-a)$ over the interval. Geometrically the
tangent at $c$ is parallel to the secant joining the endpoints. The
theorem asserts existence, not location; a function may have several
such $c$.}
\label{fig:vol02:ch05:mean-value}
\end{figure}
